\chapter{Tinea capitis}
\index{Tinea capitis}

\section*{Definition and Overview}
Tinea capitis, colloquially known as ringworm of the scalp, is a common superficial fungal infection (dermatophytosis) that primarily affects the scalp skin, hair follicles, and hair shafts. The term is derived from the Latin word ``tinea,'' meaning worm or moth, reflecting the historical belief that the lesions were caused by parasitic worms, and ``capitis,'' referring to the head. Despite its common name, tinea capitis has no relation to helminthic infections; rather, it is caused by specialized, keratinophilic fungi known as dermatophytes, which belong to the genera \textit{Trichophyton} and \textit{Microsporum}. These fungi possess the unique ability to metabolize keratin, the primary structural protein found in the outer layer of human skin, hair, and nails.

This condition is of significant clinical and public health importance. It represents one of the most common pediatric dermatoses worldwide, particularly affecting children before the age of puberty. Tinea capitis exhibits a diverse spectrum of clinical presentations, ranging from a mild, non-inflammatory scaling that mimics seborrheic dermatitis, to highly inflammatory, suppurative lesions called kerions, which can cause permanent scarring alopecia (hair loss) and systemic symptoms. Beyond the physical morbidity, tinea capitis carries a profound psychosocial impact. Affected children often face school exclusion, social isolation, and self-esteem issues due to visible hair loss and scalp lesions. Furthermore, because of its highly contagious nature, the infection can rapidly spread within families, schools, and daycare centers, necessitating prompt diagnosis, targeted systemic antifungal therapy, and robust public health measures to control transmission. Understanding the complex interplay of fungal biology, host immune response, and environmental factors is essential for effective clinical management and prevention.

\section{*Epidemiology}
Tinea capitis is a global disease, but its prevalence, incidence, and causative agents vary widely depending on geographical region, climate, socioeconomic factors, and sanitary conditions. It is predominantly a disease of childhood, most frequently observed in children aged 3 to 10 years, with a peak incidence in early school-age children. It is relatively uncommon in adults, as post-pubertal changes in sebum composition---specifically the production of medium-chain fatty acids with fungistatic properties---and increased scalp skin thickness provide a protective barrier against dermatophyte invasion. When adults are affected, they are typically postmenopausal females, individuals with severe immunodeficiency (such as HIV/AIDS or those undergoing immunosuppressive therapy), or caregivers in close contact with infected children.

The epidemiological landscape has shifted-dramatically over the past century. Historically, \textit{Microsporum audouinii} was the primary cause of epidemic tinea capitis in North America and Europe. Today, \textit{Trichophyton tonsurans} has emerged as the dominant causative organism in these regions, accounting for over 90\% of cases in the United States and the United Kingdom. This shift from \textit{Microsporum} to \textit{Trichophyton} is epidemiologically significant because \textit{T. tonsurans} infections are often non-inflammatory and subclinical, allowing asymptomatic carriers (both children and adults) to act as reservoirs, facilitating silent transmission. In contrast, in Southern and Eastern Europe, North Africa, and parts of South America, the zoophilic organism \textit{Microsporum canis}, typically acquired from domestic pets like cats and dogs, remains the primary cause. In sub-Saharan Africa, \textit{Trichophyton violaceum} and \textit{Microsporum audouinii} are highly prevalent. Socioeconomic factors, including overcrowding, low income, poor hygiene, large family sizes, and sharing of grooming utensils, are strongly associated with higher transmission rates. The disease also shows a racial predisposition in some urban areas, being disproportionately diagnosed in children of African descent, which may be linked to specific hair care practices, hair structure, or social dynamics.

\section{*Aetiology and Pathophysiology}
The development of tinea capitis requires contact with viable fungal elements, followed by adherence, penetration, and colonization of the keratinized structures of the scalp. The aetiological agents are classified into three ecological groups based on their primary host or habitat:
\begin{itemize}
    \item \textbf{Anthropophilic Dermatophytes:} These organisms are adapted exclusively to human hosts and are transmitted primarily via direct person-to-person contact or indirectly through fomites (such as combs, brushes, hats, theater seats, and pillows). Common examples include \textit{Trichophyton tonsurans}, \textit{Trichophyton violaceum}, \textit{Trichophyton schoenleinii}, and \textit{Microsporum audouinii}. These infections tend to run a chronic, mild to moderately inflammatory course, as the host immune system does not mount a vigorous inflammatory response against adapted human pathogens.
    \item \textbf{Zoophilic Dermatophytes:} These fungi are primarily pathogens of animals but can be transmitted to humans through direct contact with infected pets or farm animals, or indirectly through contaminated environments. The most notable organism is \textit{Microsporum canis} (from cats and dogs), followed by \textit{Trichophyton mentagrophytes} (from rodents) and \textit{Trichophyton verrucosum} (from cattle). Because these fungi are not adapted to human hosts, they typically elicit a robust, highly inflammatory cell-mediated immune response, often leading to inflammatory lesions like kerions.
    \item \textbf{Geophilic Dermatophytes:} These organisms reside in the soil and decompose keratinaceous debris. Infection in humans is acquired through direct contact with soil. \textit{Microsporum gypseum} is the primary geophilic species causing tinea capitis. Like zoophilic species, geophilic dermatophytes often evoke a strong inflammatory response in humans.
\end{itemize}

The pathophysiological process begins when arthroconidia (fungal spores) adhere to the stratum corneum of the scalp. Germination of these spores allows hyphae to penetrate the stratum corneum and grow downward into the hair follicle, utilizing keratinases, proteolytic enzymes, and lipolysins to digest the keratin. Once inside the follicle, the fungus invades the hair shaft. Depending on the species, the infection of the hair shaft follows one of three distinct patterns:
\begin{itemize}
    \item \textbf{Ectothrix Invasion:} Arthroconidia form a sheath around the outside of the hair shaft, and the cuticle of the hair is destroyed. Hyphae also grow inside the hair shaft. Under Wood's lamp examination, ectothrix infections caused by certain \textit{Microsporum} species (such as \textit{M. canis} and \textit{M. audouinii}) emit a bright blue-green fluorescence due to the metabolic production of pteridine. Ectothrix pathogens include \textit{Microsporum canis}, \textit{Microsporum audouinii}, and \textit{Trichophyton verrucosum}.
    \item \textbf{Endothrix Invasion:} Fungal growth and arthroconidia formation are confined entirely within the hair shaft, leaving the outer cuticle intact. Hairs infected in an endothrix pattern are extremely fragile and break off at the level of the scalp surface. Endothrix infections do not fluoresce under Wood's lamp. The primary endothrix pathogens are \textit{Trichophyton tonsurans} and \textit{Trichophyton violaceum}.
    \item \textbf{Favus:} This is a unique and severe form of infection caused almost exclusively by \textit{Trichophyton schoenleinii}. Hyphae grow longitudinally inside the hair shaft, and air bubbles (tunnels) are formed within the hair structure, rather than dense conidia. This results in the characteristic clinical presentation of scutula and scarring.
\end{itemize}

\section*{Clinical Features}
The clinical presentation of tinea capitis is highly variable and depends on the causative organism (anthropophilic versus zoophilic), the type of hair invasion (endothrix versus ectothrix), and the intensity of the host's cell-mediated immune response. The primary clinical classifications include:
\begin{itemize}
    \item \textbf{Non-Inflammatory Tinea Capitis (Grey Patch Type):} Typically caused by ectothrix organisms like \textit{Microsporum canis} or \textit{Microsporum audouinii}. It begins as small, erythematous papules around the hair follicle that expand centrifugally to form circular, scaly, greyish patches. The hair in the affected areas becomes dull, lusterless, and brittle, breaking off a few millimeters above the scalp surface. This results in partial alopecia with stubby, ``grey'' hairs. Pruritus (itching) is common, and the lesions are often multiple and well-demarcated.
    \item \textbf{Non-Inflammatory Tinea Capitis (Black Dot Type):} Characteristically caused by endothrix organisms like \textit{Trichophyton tonsurans} or \textit{Trichophyton violaceum}. Because the hair shaft is structurally weakened from within, the hair breaks off precisely at the level of the scalp skin. This leaves behind small, dark, hyperkeratotic remains of the hair within the follicular openings, presenting as multiple ``black dots'' scattered across areas of diffuse scaling and patchy alopecia. This type can easily be mistaken for seborrheic dermatitis, as erythema may be minimal.
    \item \textbf{Inflammatory Tinea Capitis (Kerion):} A kerion (or kerion celsi) represents a severe, T-cell-mediated hypersensitivity reaction to the fungus, most frequently associated with zoophilic species like \textit{Microsporum canis} or \textit{Trichophyton verrucosum}. It presents as a painful, boggy, highly inflammatory, fluctuant mass studded with pustules. The surface may drain pus from multiple follicular openings, resembling a bacterial carbuncle or abscess. The hair in the affected area falls out easily or can be painlessly plucked. Kerions are frequently accompanied by regional lymphadenopathy (particularly occipital, posterior auricular, or cervical lymph nodes), fever, malaise, and a generalized cutaneous eruption known as a dermatophytid (id) reaction. If not treated promptly, kerions result in irreversible damage to hair follicles, culminating in scarring alopecia and permanent hair loss.
    \item \textbf{Favus (Tinea Favosa):} A rare, chronic form of tinea capitis caused by \textit{Trichophyton schoenleinii}. It is characterized by the formation of sulfur-yellow, cup-shaped crusts called scutula that surround loose, loop-like hair shafts. These crusts consist of dense masses of mycelia, epithelial debris, and inflammatory cells. Over time, these lesions coalesce, emitting a characteristic ``mousy'' or ``musty'' odor. Favus leads to extensive scalp atrophy, scarring, and permanent alopecia.
\end{itemize}

\section*{Differential Diagnosis}
Because the clinical features of tinea capitis can mimic several other dermatological conditions of the scalp, a thorough differential diagnosis is essential. The primary conditions to consider include:
\begin{itemize}
    \item \textbf{Seborrheic Dermatitis:} Often presents with diffuse scaling and mild erythema, closely resembling the seborrheic-like variant of tinea capitis. However, seborrheic dermatitis typically affects the eyebrows and nasolabial folds as well, lacks the characteristic hair fragility, broken hairs, or ``black dots,'' and does not cause patchy alopecia.
    \item \textbf{Scalp Psoriasis:} Characterized by well-demarcated, erythematous plaques covered with thick, silvery-white scales. While psoriasis can cause temporary hair shedding (telogen effluvium), it does not cause the structural hair breakage or localized patches of alopecia characteristic of tinea capitis. Wood's lamp examination and fungal cultures are negative.
    \item \textbf{Alopecia Areata:} An autoimmune condition presenting as smooth, completely hairless patches without signs of inflammation, scaling, or erythema. The pathognomonic ``exclamation point'' hairs may be seen at the margins of active patches, but there are no scaly plaques or pustules, and fungal investigations are negative.
    \item \textbf{Trichotillomania:} A psychiatric disorder characterized by the compulsive urge to pull out one's own hair. It presents as irregular patches of hair loss with hairs of varying lengths. Unlike tinea capitis, the scalp skin appears normal without scaling, erythema, pustules, or black dots (except for normal follicular plugs).
    \item \textbf{Bacterial Folliculitis and Furunculosis:} Acute bacterial infections (often caused by \textit{Staphylococcus aureus}) that present as painful follicular pustules or deep nodules. While they can mimic the pustular appearance of a kerion, they lack the underlying fungal scaling, characteristic hair shaft fragility, and do not respond to antifungal therapy. A bacterial culture will help differentiate the two.
    \item \textbf{Lichen Planopilaris and Discoid Lupus Erythematosus:} Chronic inflammatory conditions that cause scarring (cicatricial) alopecia. Lichen planopilaris presents with perifollicular erythema and scaling, while discoid lupus shows active erythematous plaques with follicular plugging, atrophy, and telangiectasia. These conditions typically affect adults rather than children and are characterized by histographical findings of lichenoid or vacuolar interface dermatitis, with negative fungal cultures.
    \item \textbf{Atopic Dermatitis:} Can involve the scalp, causing scaling, crusting, and intense pruritus. However, it is usually associated with eczematous lesions on other classic body sites (flexural creases) and a personal or family history of atopy, and it does not feature localized hair breakage or fungal elements.
\end{itemize}

\section*{When to Seek Medical Care}
Scalp changes in children should always be evaluated by a healthcare professional. Parents or caregivers should schedule an appointment with a general practitioner or dermatologist if a child develops localized hair loss, persistent scaling of the scalp, itching, or unusual red patches on the head. Early intervention is critical to prevent the infection from progressing to inflammatory stages.

Immediate medical attention is required if the child exhibits signs of a severe inflammatory response or secondary bacterial infection. Caregivers should seek prompt medical evaluation if the scalp develops a painful, swollen, boggy, or pus-filled lump (suggestive of a kerion), or if the infection is accompanied by a high fever, rapidly spreading redness, severe pain, or swelling of the face or neck.

While tinea capitis itself is not a medical emergency, severe secondary bacterial infections (such as cellulitis) can occasionally lead to systemic complications if left untreated. In the rare event that a child exhibits life-threatening symptoms---such as difficulty breathing, altered consciousness, a stiff neck, or extreme lethargy---caregivers must call emergency services immediately (for example, local emergency services in many countries, 999 in the United Kingdom, or 911 in the United States) or present to the nearest emergency department.

\section*{Diagnosis and Investigations}
Confirming the diagnosis of tinea capitis through objective testing is imperative before initiating treatment, as systemic antifungal therapies carry potential risks and require long courses. The diagnostic workup involves several clinical and laboratory tools:
\begin{itemize}
    \item \textbf{Wood's Lamp Examination:} A diagnostic tool that utilizes long-wave ultraviolet light (365 nm). When directed at the scalp in a darkened room, ectothrix infections caused by \textit{Microsporum canis} or \textit{Microsporum audouinii} exhibit a bright, pathognomonic blue-green fluorescence. This fluorescence is caused by a metabolite (pteridine) produced by the active fungus. However, endothrix infections (such as those caused by \textit{Trichophyton tonsurans}) do not fluoresce, meaning a negative Wood's lamp test does not rule out tinea capitis.
    \item \textbf{Dermoscopy (Trichoscopy):} A non-invasive, in-office technique using a dermatoscope to visualize the scalp and hair shafts at high magnification. Specific trichoscopic markers for tinea capitis include ``comma hairs'' (short, bent hairs), ``corkscrew hairs'' (coiled hairs), ``bar-code hairs'' (horizontal white bands), broken hairs, and Morse code-like hairs. Trichoscopy is highly sensitive and aids in differentiating tinea capitis from alopecia areata or trichotillomania.
    \item \textbf{Direct Microscopy (KOH Prep):} A rapid and inexpensive test. Hair shafts and skin scrapings are collected from the active border of the lesion (or by using a clean toothbrush to brush the scalp) and placed on a glass slide. A few drops of 10\% to 20\% potassium hydroxide (KOH) solution are added, sometimes with a fungal stain (like chlorazol black E or Parker ink). The slide is gently heated to dissolve the human keratin, leaving the fungal structures intact. Under the microscope, the clinician can visualize branching, septate hyphae and arthroconidia, confirming a fungal infection and identifying whether it is an ectothrix or endothrix pattern.
    \item \textbf{Fungal Culture:} The historical gold standard for diagnosis. Samples are inoculated onto Sabouraud dextrose agar (often containing cycloheximide and chloramphenicol to inhibit contaminating molds and bacteria) or Dermatophyte Test Medium (DTM). The culture is incubated at room temperature for 2 to 4 weeks. Culture allows for the definitive identification of the specific fungal species, which is crucial for determining the source of infection (human versus animal) and selecting the most effective antifungal agent.
    \item \textbf{Polymerase Chain Reaction (PCR) Assays:} Modern molecular diagnostic techniques that detect dermatophyte DNA directly from hair and skin samples. PCR assays are highly sensitive and specific, providing results within 24 to 48 hours, significantly faster than fungal cultures. Although more expensive and less widely available in routine clinical practice, PCR is increasingly utilized in reference laboratories.
\end{itemize}

\section*{Management}
Topical antifungal creams, lotions, and shampoos are insufficient as monotherapy for tinea capitis because they cannot penetrate the deep keratinized structures of the hair follicle root. Therefore, systemic oral antifungal therapy is mandatory for a clinical and mycological cure. The selection of the agent and dosage depends on the causative fungal species, the patient's age, and safety profiles:
\begin{itemize}
    \item \textbf{Griseofulvin:} For decades, griseofulvin was the gold standard and remains highly effective, particularly for infections caused by \textit{Microsporum} species (such as \textit{M. canis}). It is fungistatic and works by disrupting the mitotic spindle structure of fungal cells. The typical pediatric dosage is 20 to 25 mg/kg/day of microsize formulation (or 10 to 15 mg/kg/day of ultramicrosize formulation) administered daily for 6 to 8 weeks, or up to 12 weeks. Griseofulvin must be taken with a fatty meal (e.g., milk, ice cream, or peanut butter) to enhance systemic absorption. Potential side effects include headaches, gastrointestinal upset, photosensitivity, and rare hepatotoxicity.
    \item \textbf{Terbinafine:} An allylamine antifungal that is fungicidal, inhibiting squalene epoxidase, an essential enzyme in fungal cell membrane synthesis. Terbinafine is highly effective against \textit{Trichophyton} species (specifically \textit{T. tonsurans}) and has become the first-line treatment in many regions due to its shorter treatment duration. The standard treatment course is 2 to 4 weeks, with pediatric dosing based on weight bands: 62.5 mg/day for children weighing less than 20 kg, 125 mg/day for 20--40 kg, and 250 mg/day for those over 40 kg. It is less effective against \textit{Microsporum} species, which may require higher doses and longer courses (6--8 weeks). Side effects are generally mild, including gastrointestinal disturbances, taste alterations, rash, and transient liver enzyme elevations.
    \item \textbf{Itraconazole:} A triazole antifungal that inhibits lanosterol 14-alpha-demethylase, halting ergosterol synthesis. It is effective against both \textit{Trichophyton} and \textit{Microsporum} species. It can be administered as a continuous daily therapy (3 to 5 mg/kg/day) for 4 to 6 weeks, or as pulse therapy (5 mg/kg/day for 1 week per month, repeated for 2 to 3 consecutive months). It is available as an oral solution, making it easier to administer to very young children.
    \item \textbf{Fluconazole:} Another triazole antifungal with excellent absorption and a favorable safety profile. It is administered daily at a dose of 6 mg/kg/day for 3 to 6 weeks. Like itraconazole, it is available in a liquid suspension and is well-tolerated by pediatric patients.
    \item \textbf{Adjunctive Topical Therapies:} While topical treatments cannot cure the infection alone, adjunctive use of antifungal shampoos is strongly recommended to reduce the shedding of viable fungal spores, thereby lowering infectivity and preventing transmission to others. Patients and close contacts should use selenium sulfide (1\% or 2.5\%), ketoconazole (2\%), or ciclopirox shampoo 2 to 3 times per week. The shampoo should be left on the scalp for 5 to 10 minutes before rinsing.
    \item \textbf{Corticosteroids:} In cases of severe, highly inflammatory kerion celsi, oral corticosteroids (e.g., prednisolone at 1 mg/kg/day for the first 1 to 2 weeks of antifungal therapy) are sometimes prescribed to rapidly suppress the host immune response. The goal is to alleviate severe pain, reduce swelling, and decrease the risk of permanent scarring alopecia, although clinical trials have shown mixed results regarding their long-term efficacy in preventing scarring.
\end{itemize}

\section*{Complications}
If tinea capitis is left untreated, misdiagnosed, or treated with inadequate doses or durations of antifungal medication, several complications can arise:
\begin{itemize}
    \item \textbf{Kerion Formation:} A progression from a simple non-inflammatory infection to a highly inflammatory, purulent mass. As detailed previously, this is a hypersensitivity reaction that can cause significant pain, systemic symptoms (fever, lymphadenopathy), and tissue destruction.
    \item \textbf{Scarring Alopecia (Cicatricial Alopecia):} The most significant long-term complication of severe inflammatory tinea capitis (kerion) or favus. The intense inflammatory response destroys the stem cells located in the hair follicle bulge, preventing future hair regeneration. This results in permanent patches of baldness and fibrotic scarring of the scalp, which can have lifelong psychological consequences for the patient.
    \item \textbf{Secondary Bacterial Infection (Impetiginization):} The inflamed, broken scalp skin in tinea capitis, particularly in kerion or severely pruritic lesions, provides an entry portal for pathogenic bacteria, most commonly \textit{Staphylococcus aureus} or \textit{Streptococcus pyogenes}. This can result in secondary impetigo, cellulitis, or localized abscess formation, requiring concurrent oral antibiotic therapy.
    \item \textbf{Dermatophytid (Id) Reaction:} An immunological hypersensitivity reaction to fungal antigens that occurs in up to 5\% of patients, often shortly after starting systemic antifungal therapy. It manifests as a symmetric, highly pruritic, papular, follicular, or vesicular eruption, typically starting on the face and neck before spreading to the trunk and extremities. It does not contain active fungal elements and resolves spontaneously as the primary infection is cleared; it should not be mistaken for an allergic drug reaction.
    \item \textbf{Psychosocial Impairment:} The presence of visible hair loss, scaly patches, or inflammatory masses on the head can lead to significant psychological distress, anxiety, depression, low self-esteem, and social withdrawal in school-aged children. School exclusions, though often unnecessary if appropriate treatment has begun, compound this impact.
\end{itemize}

\section*{Prognosis and Outcomes}
The prognosis for individuals diagnosed with tinea capitis is generally excellent, provided that appropriate systemic antifungal therapy is initiated early and completed as prescribed. Mycological and clinical cure rates exceed 80\% to 90\% with modern first-line agents like terbinafine and griseofulvin.

For the non-inflammatory variants (grey patch and black dot), complete hair regrowth is the standard outcome. Regrowth begins slowly and may take several months after the completion of antifungal therapy to become fully visible. The texture and color of the new hair may temporarily differ from the surrounding hair but typically normalizes over time.

In contrast, the prognosis for hair regrowth is guarded in patients who develop kerions or favus. If the inflammatory process is severe or prolonged, the destruction of hair follicles leads to permanent scarring alopecia. The size of the permanent bald patch correlates directly with the delay in seeking medical care and starting systemic treatment. Once scarring has occurred, medical treatments cannot restore hair, and patients may require surgical options, such as hair transplantation or scalp reduction, in adulthood. Additionally, treatment failure or relapse can occur if there is poor compliance (due to the long course or unpalatable taste of medications), inadequate absorption (e.g., taking griseofulvin without fatty foods), or continuous re-exposure to asymptomatic carriers or infected pets in the household.

\section*{Prevention and Self-Care}
Preventing the acquisition and transmission of tinea capitis requires a combination of personal hygiene, environmental disinfection, and vigilance. Key preventative measures include:
\begin{itemize}
    \item \textbf{Avoid Sharing Personal Items:} Children should be educated never to share items that come into contact with the scalp, including combs, hairbrushes, hair ties, hats, caps, helmets, hoods, pillows, and towels. Fungal spores can survive on these fomites for months and serve as a source of infection.
    \item \textbf{Regular Disinfection of Grooming Tools:} Combs, brushes, and hair accessories used by an infected individual should be disinfected daily by soaking them in a solution of chlorine bleach diluted with water (1:10 ratio) or an appropriate antifungal disinfectant for at least 10 minutes, or they should be discarded.
    \item \textbf{Household Contact Screening:} Because asymptomatic carriage is common, all family members and close contacts of an infected individual should be screened for signs of infection. Using an antifungal shampoo (e.g., ketoconazole 2\%) twice weekly for all household members during the treatment period of the index case is a prudent measure to prevent reinfection.
    \item \textbf{Veterinary Evaluation of Pets:} Since zoophilic dermatophytes (especially \textit{Microsporum canis}) are highly contagious, domestic pets (cats, dogs, rabbits) in the household must be evaluated by a veterinarian, even if they do not display obvious signs of hair loss or scaling. Asymptomatic feline carriers are a frequent cause of treatment failure in families.
    \item \textbf{Good Hand Hygiene:} Wash hands thoroughly with soap and warm water after handling pets, playing outdoors, or touching soil.
    \item \textbf{School and Daycare Communication:} Inform the child's school or daycare center of the diagnosis so that staff can monitor other children for symptoms and implement enhanced cleaning protocols. Children may usually return to school once they have commenced appropriate oral antifungal therapy and adjunctive shampooing.
    \item \textbf{Proper Hair and Scalp Hygiene:} Wash the scalp regularly. After hair washes, ensure the hair and scalp are dried completely, as fungi thrive in warm, moist environments. Avoid tight hairstyles that trap moisture or cause mechanical trauma to the scalp.
\end{itemize}

\section*{Sources and Further Reading}
\begin{itemize}
    \item \textbf{World Health Organization (WHO):} Guidelines on the management of common skin diseases in children. URL: \texttt{https://www.who.int}
    \item \textbf{American Academy of Dermatology (AAD):} Comprehensive patient education resources on ringworm of the scalp. URL: \texttt{https://www.aad.org}
    \item \textbf{British Association of Dermatologists (BAD):} Clinical guidelines and patient information leaflets regarding tinea capitis. URL: \texttt{https://www.bad.org.uk}
    \item \textbf{Centers for Disease Control and Prevention (CDC):} Fungal Diseases Branch resources on dermatophyte infections and ringworm prevention. URL: \texttt{https://www.cdc.gov/fungal/diseases/ringworm/}
    \item \textbf{DermNet NZ:} Authoritative, peer-reviewed clinical information and photographic library of tinea capitis and kerions. URL: \texttt{https://dermnetnz.org/topics/tinea-capitis}
    \item \textbf{National Health Service (NHS) UK:} General public health advice on ringworm symptoms, treatments, and school exclusion policies. URL: \texttt{https://www.nhs.uk/conditions/ringworm/}
    \item \textbf{Royal Children's Hospital (RCH) Melbourne:} Clinical practice guidelines on ringworm (tinea) and its management in pediatric settings. URL: \texttt{https://www.rch.org.au/clinicalguide/guideline\_index/Ringworm/}
\end{itemize}